% Transcription par lecture visuelle directe, sans OCR.
% Statut : transcrit ; relecture indépendante à effectuer.
% Chapitre complet, p. 331 (suite)–349 (début) ; le chapitre troisième commence ensuite p. 349.
% Source : PDF original, vue 165, côté droit, page manuscrite 331 (suite).
\chapter{Suite de l'organisation de la Congrégation des hommes. Œuvres de ses sections. Nouvelles œuvres de celle des jeunes gens.}

La fondation de la Congrégation des hommes était un bien grand événement pour la Congrégation considérée dans son ensemble ; elle doublait ses forces, multipliait ses rapports et agrandissait son action. Voilà aujourd'hui deux sociétés bien distinctes, qui vont unir leurs efforts pour parvenir au même but, et cependant ces deux sociétés n'en forment réellement qu'une car leurs règlements expriment formellement le vœu qu'elles se considèrent comme ne faisant qu'une seule famille ; c'est par cette raison qu'à l'avenir si nous sommes obligés de distinguer ces deux branches d'une même société par les dénominations de Congrégation des jeunes gens et de Congrégation des hommes lorsque nous aurons à rendre compte des faits qui les concernent personnellement, lorsque nous aurons à rappeler le souvenir des affaires générales auxquelles elles ont pris part collectivement, nous n'en continuerons pas moins à dire purement et simplement la Congrégation.

Nous ne pourrons plus maintenant, ainsi que nous l'avons déjà annoncé, vous signaler dans ces mémoires, d'année en année l'époque du gouvernement de chaque préfet, comme nous l'avons fait dans l'histoire de la Congrégation des jeunes gens. Pour
% Source : PDF original, vue 166, côté gauche, page manuscrite 332.
y suppléer autant qu'il est en nous, nous allons chercher du moins à vous faire connaître, avant d'aller plus loin, ceux qui jusqu'à ce jour ont successivement présidé l'une et l'autre Congrégation. Si j'ai oublié quelques noms, vous excuserez l'infidélité de ma mémoire et vous rectifierez cette omission au moyen des documents que vous avez entre les mains.

\section*{Liste des Préfets des jeunes gens}
% Mention marginale de la page manuscrite 332, placée ici avant la liste.
Dans la Congrégation des jeunes gens, M. André Terret succéda à M. Matthieu Garnier, sous la préfecture duquel s'était opérée la division ; ensuite ont été successivement Préfets MM. Desgorges, Delorme, de Villers, Empaire aîné, Montetlan, Guérin, Chalandon et enfin M. Desplace préfet actuel. Parmi ces chefs que la Congrégation des jeunes gens a placés tour à tour à sa tête, il en est deux, MM. Delorme et de Villers que notre commun maître a déjà appelés à lui pour recevoir la récompense après laquelle nous soupirons tous.

M. Delorme, quoique encore bien jeune, avait débuté avec distinction dans la carrière du barreau. Il était rempli de talent et de modestie. Sa grande piété avait marqué sa place dans la Congrégation ; son mérite y fut bien vite apprécié. Ses vertus firent l'édification de ses amis. La rectitude de son jugement le proposait au gouvernement de la Congrégation. Elle s'estima donc heureuse de le choisir pour préfet ; mais il ne lui fut pas donné de rendre en cette qualité de bien longs services à ses amis. Ce fut pendant qu'il était en exercice qu'il fut atteint de la maladie grave qui en peu de jours l'a conduit au tombeau. C'est ainsi qu'à la fleur de son âge il fut mis en possession de cette récompense éternelle dont la bonté de notre Dieu lui avait fait connaître tout le prix dès sa plus
% Source : PDF original, vue 166, côté droit, page manuscrite 333.
tendre jeunesse.

M. de Villers, littérateur distingué a servi pendant toute sa vie le Seigneur avec fidélité. Ses dernières années surtout ont été entièrement consacrées à la défense de la religion. Non seulement il remplissait avec exactitude ses devoirs de Congréganiste, mais encore il faisait usage de ses talents pour plaider la cause de l'éternelle vérité et la défendre victorieusement avec toutes les armes que lui fournissait une logique serrée contre les attaques sans cesse renouvelées auxquelles elle est en butte dans notre continent ; et de plus il traversait pour ainsi dire les mers, à l'aide de l'œuvre de la Propagation de la foi, pour la faire triompher aussi dans le nouveau monde. Il fut un des fondateurs et un des plus fermes appuis de cette œuvre admirable. Il en a pendant quelque temps rédigé les annales.

M. de Villers avait rempli dans le temps avec beaucoup de zèle et toute la supériorité de ses qualités personnelles les fonctions de Préfet de la Congrégation des jeunes gens. Il possédait dans un corps bien débile une âme bien forte. Si comme je l'ai souvent entendu dire, le style est tout l'homme, à en juger par l'élévation du sien, ce digne et vertueux ami s'était élevé à une si grande hauteur que les esprits vulgaires auraient eu de la peine à l'y suivre.

\section*{Préfets des hommes}
% Mention marginale de la page manuscrite 333 ; date marginale : 1817–1820.
La présidence de la Congrégation des hommes a été confiée successivement à MM. Jacquier, Victor de Verna, Bonnet, Terret et à l'auteur de ces mémoires. Quelques-uns d'entre eux y ont été appelés plusieurs fois.

En me disposant à transcrire ici ces noms, je vous avouerai franchement ma foiblesse. Je regrettais de m'être imposé
% Source : PDF original, vue 167, côté gauche, page manuscrite 334.
rigoureusement la loi de ne me permettre dans cet écrit aucun éloge de mes amis, à moins que la mort ou une absence prolongée ne les eût éloignés de la Congrégation. Il m'était pénible de ne pas pouvoir par ce motif satisfaire l'élan de mon cœur en disant à mes frères tout ce que je pense de celui d'entre nous qu'il me semblait que je pouvais leur présenter comme le plus parfait modèle du véritable Congréganiste. J'en étais vivement affecté et même troublé, et voilà que la Providence vient elle-même m'affranchir, bien douloureusement à la vérité, de ces entraves en faisant parvenir jusqu'à moi au moment précis où j'hésitais encore la triste nouvelle de la mort de ce parfait ami. Voilà donc que j'ai acquis d'une manière sans doute bien cruelle le droit de m'exprimer en pleine liberté sur ce qui concerne l'ami que nous pleurons, M. Victor de Verna\footnote{La marge porte une croix et « 17 juin 1841 ». Date lue sur le scan, à distinguer de la période du récit. (Note du transcripteur).}.

\section*{M. de Verna}
% Mention marginale des pages manuscrites 334 et 336.
Je ne m'arrêterai cependant ni à la fermeté d'âme ni au beau caractère qu'il a déployé comme homme politique. Ainsi que moi, vous l'avez connu et apprécié sous ce rapport. D'autres vraisemblablement ne manqueront pas de rappeler sa noble conduite dans les circonstances délicates où il s'est trouvé. Pour moi je ne veux le considérer que comme Congréganiste. C'est là seulement ce qui rentre dans le but que je me suis proposé.

M. de Verna avait fait partie de la Congrégation de la Ste Famille. Dès que la nouvelle Congrégation des hommes fut établie, nous nous empressâmes de le rechercher ; il nous combla de joie en consentant à s'unir à nous. Dès son début il se fit remarquer par son assiduité à tous les exercices, sa soumission entière à tout ce que les chefs de la Congrégation
% Source : PDF original, vue 167, côté droit, page manuscrite 335.
exigeaient de lui, son zèle pour la gloire de Dieu, sa grande charité envers le prochain, sa bonté et sa douceur dans tous ses rapports avec les Congréganistes. Ce mérite éminent ne pouvait pas rester inaperçu. Il fut bien vite apprécié de tous, et notre digne ami fut d'abord chargé de la présidence d'une section, et bientôt après appelé au gouvernement de la Congrégation premièrement comme assistant et ensuite comme Préfet. Plus ces divers emplois plaçaient les Congréganistes à portée de l'observer de près, plus aussi ils trouvaient des occasions d'admirer de plus en plus son excellent jugement.

M. de Verna plus que personne avait toujours en vue dans tout ce qu'il entreprenait ce qui pouvait procurer la gloire de Dieu. Aussi a-t-il pris une part active à toutes les œuvres de zèle auxquelles soit directement soit indirectement la Congrégation a été appelée à concourir.

Au milieu de toutes ces œuvres qu'il serait trop long d'énumérer ici, nous distinguerons cependant celle de la Propagation de la Foi, dans laquelle il a présidé constamment jusqu'à la mort le conseil central de Lyon. Vous le savez tout comme moi, Messieurs, notre vénérable ami était devenu si parfaitement à Lyon l'homme par excellence, l'homme nécessaire, qu'il semblait qu'aucune œuvre ne pouvait avoir de succès s'il n'en prenait en main la direction. Il était pour notre ville ce que M. de Montmorency avait été pour Paris. Comme lui on aurait bien pu l'appeler aussi le ministre des bonnes œuvres.

Qu'il me soit permis d'ajouter à ce détail des vertus dont il vous a rendus le témoin, une seule réflexion. Comment cette vertu de M. de Verna est-elle parvenue à une si grande hauteur,
% Source : PDF original, vue 168, côté gauche, page manuscrite 336.
comment s'est-elle élevée comme un arbre vigoureux, étendant ses exemples comme un épais feuillage sur notre Congrégation pour y conserver sous leur ombrage tutélaire l'esprit de zèle et de ferveur ? Voilà tout le secret, Messieurs ; c'est que avant de s'élever ainsi cet arbre avait poussé de profondes racines aux pieds de notre souverain maître par la pratique d'une humilité sincère, et qu'une oraison assidue faisait couler par torrent sur cette belle âme les eaux vives de la grâce. Ne soyons pas étonnés qu'avec de semblables moyens nous ayons vu cette éminente vertu prendre un accroissement aussi rapide, car c'est ainsi que Dieu sanctifie ses élus.

Remercions donc le ciel, remercions Marie notre mère de nous avoir associés à cette âme privilégiée. Que son souvenir se perpétue parmi nous d'âge en âge, qu'il y demeure et qu'il devienne pour la Congrégation comme une précieuse semence qui lui produise des imitateurs. Si la Providence a bien voulu nous montrer un semblable modèle, c'est pour nous engager à y conformer notre conduite. Si dans notre foiblesse nous éprouvons quelques difficultés à nous en rapprocher, ah ! du moins efforçons-nous de le suivre de loin.

Si la mémoire de ces saints amis dont les exemples nous ont édifiés pendant leur passage sur cette misérable terre doit être pour nous un objet de vénération, il me semble aussi que rien ne peut autant redoubler notre confiance ; car si pendant leur vie ils ont eu pour nous une charité si expansive, jugez, si vous en êtes capables, combien cet amour doit être devenu plus vif maintenant que leurs sentiments pour nous se sont épurés, réchauffés et en quelque sorte déifiés en se retrempant dans la
% Source : PDF original, vue 168, côté droit, page manuscrite 337.
charité de Dieu lui-même auquel ils sont inséparablement unis pour l'éternité. Voilà pourtant les intercesseurs que nous avons trouvés ; ne sommes-nous pas bien fondés à espérer qu'ils feront descendre sur notre chère Congrégation et sur nous tous, ces grâces abondantes qui nous sont si nécessaires pour soutenir notre courage au milieu des maux divers que nous avons à supporter en traversant le désert de cette vie. Ces grâces se répandront sur nous comme une douce rosée ; nous nous empresserons de recueillir cette manne céleste ; nous y trouverons une nourriture délicieuse qui conservera la vie à notre âme. Jusqu'au moment marqué par la Providence où à notre tour, nous verrons enfin s'abaisser pour nous la barrière, qui nous sépare encore de la Terre Promise, l'objet de tous nos désirs. \textit{O quando sancta se dabit quæ nescit hostem patria !}\footnote{Littéralement : « Ô quand se donnera-t-elle, la sainte patrie qui ne connaît pas d'ennemi ! » Formule attestée dans l'hymne \textit{Grates, peracto jam die}, aux complies. (Note du transcripteur).}

Vous allez peut-être, Messieurs, trouver que votre historien est un peu prolixe. Pardon, mais je n'ai pas eu le courage de résister lorsque je trouve l'occasion de m'entretenir avec vous, de ce qui fait l'objet de toutes nos pensées et de toutes nos espérances. Nous allons pourtant en revenir à la Congrégation des hommes. Nous l'avons laissée à la fin du chapitre précédent au moment où elle venait d'être fondée. Ce qu'elle avait alors de plus urgent était de se recruter ; elle eut pu sans doute le faire très rapidement, si elle eût admis de suite tous les membres de la Congrégation des jeunes gens qui par leur âge étaient appelés à entrer dans la nouvelle Congrégation, mais pour ne pas affoiblir la Congrégation mère il fallait opérer ce transport avec beaucoup de lenteur afin de lui laisser le temps de réparer par de nouvelles réceptions les enlèvements
% Source : PDF original, vue 169, côté gauche, page manuscrite 338.
qu'on était forcé de lui faire. La Congrégation des hommes devait donc chercher à se recruter en dehors ; c'était là, la première opération dont tout naturellement elle devait s'occuper ; aussi tous les travaux de ses membres se bornèrent pour ce premier moment à former tous ensemble, sous la présidence de leur préfet, le comité du zèle. Il ne fut pas difficile de trouver des noms en abondance pour remplir les listes de ce comité ; il se rencontrait tout autour de la Congrégation un grand nombre de fervents chrétiens auxquels leur âge ou leur état n'avait jamais permis à la Congrégation des jeunes gens de penser ; d'autres qu'on aurait reçus depuis quelques années dans les circonstances qui avaient suspendu les réceptions, et que dans l'intervalle ou s'étaient mariés, ou étaient devenus trop âgés ; enfin un grand nombre d'anciens membres de la Congrégation de la Ste famille possédait à un haut degré toutes les qualités que nous pouvions désirer dans nos membres ; ils étaient déjà pour ainsi dire façonnés à la Congrégation, il nous convenait bien de les rattacher à nous. Il y avait quelques années que la Congrégation de la Ste Famille s'était refusée à faire disparaître entre elle et nous la barrière du scrutin, et c'était précisément aujourd'hui les mêmes Congréganistes qu'on avait voulu se réserver de soumettre à cette formalité, qui y soumettaient à leur tour les restes épars de cette Congrégation.

\section*{Réceptions dans la Congrégation des hommes}
% Mention marginale de la page manuscrite 338, placée ici avant le bilan des réceptions.
Avec un si grand nombre de matériaux tout préparés les réceptions purent se succéder avec rapidité. La première, vu l'urgence, eut lieu le jour de la fête de la Toussaint, et quoique cette fête ne dût pas à l'avenir être un jour de réception, une
% Source : PDF original, vue 169, côté droit, page manuscrite 339.
seconde réception vint encore augmenter la Congrégation au jour de la fête de la Conception ; elle se trouva alors composée de dix-huit à vingt membres. On crut ce nombre suffisant pour compléter l'organisation, en nommant un trésorier, un secrétaire et en formant les trois sections ; ce n'était encore pour ainsi dire que des cadres, mais on commençait toujours à agir en attendant de pouvoir faire beaucoup. Les assemblées générales se tenaient avec exactitude et complétaient jusques dans les plus petits détails par des arrêtés, tout ce qui était nécessaire, pour achever d'organiser la nouvelle société. Ainsi on choisit l'emblème qui devait servir de sceau à la Congrégation. Celle des jeunes gens avait adopté le chiffre de sa Mère, un M et un A entrelacés et couronnés. Celle des hommes choisit le cœur de Marie. Elle ne crut pas pouvoir mieux se placer. Elle y ajouta pour légende ces mots : \textit{Sine macula}\footnote{Littéralement : « Sans tache ». (Note du transcripteur).}. Ce rapprochement réunissait les deux titres de la Congrégation et proclamait de nouveau le glorieux privilège de l'Immaculée Conception.

La Congrégation des hommes fut encore appelée à délibérer sur le choix des patrons de ses sections et comités. St François Xavier fut choisi pour le comité du zèle. Depuis des siècles St Bonaventure était considéré comme le patron des enfants de Lyon ; c'était donc lui confier le soin de sa famille que de le choisir pour le patron de la section de l'éducation. St Martin, le héros de la charité chrétienne, devint le patron de la section de charité ; et on pensa que pour guider la section de l'instruction dans le dessein qu'elle avait conçu de faire briller les lumières de la foi aux yeux de tant d'aveugles, on ne pouvait pas mieux s'adresser qu'à notre apôtre St Irénée de qui nous les avons reçues.

% Source : PDF original, vue 170, côté gauche, page manuscrite 340.
\section*{Sections des hommes}
% Mention marginale de la page manuscrite 340.
La Congrégation des hommes s'augmentait de plus en plus. Chaque nouvelle fête amenait de nouveaux membres ; aussi dès la fin de la première année, elle se composait déjà de 40 Congréganistes, et au bout de deux ou trois ans, elle en avait de 60 à 70. Il ne fut pas difficile aux sections d'entreprendre les œuvres dont elles étaient chargées. Un assez grand nombre d'anciens Congréganistes se trouvant en faire partie, il leur fut facile de mettre promptement au courant de nos usages ceux qui débutaient. Aussi la manière de tenir les séances, celle de diriger les œuvres, les coutumes consacrées par le temps, tout fut exactement copié de la Congrégation des jeunes gens. C'était un arbre nouveau greffé sur un ancien ; il devait porter des fruits de même nature.

La section de l'éducation n'eut besoin que d'un travail de quelques mois pour devenir la tutrice de quinze à vingt jeunes gens qu'elle avait mis en apprentissage ; celle de charité devint une seconde section des aumônes ; elle eut bientôt rassemblé autour d'elle 30 à 40 familles honteuses. On prit des précautions pour éviter les doubles emplois et l'on eut soin de confier de préférence aux hommes celles des familles qui paraissaient plus convenablement placées entre leurs mains. Les œuvres du ressort de la section de l'instruction semblaient devoir donner des résultats beaucoup plus intéressants, mais aussi leur exercice offrait plus de difficultés. Cependant il ne lui fallut pas un laps de temps bien considérable pour se voir en état d'occuper suffisamment tous ses membres. Ce furent dès le principe les unions illicites à régulariser qui offrirent le plus d'aliment au zèle des membres
% Source : PDF original, vue 170, côté droit, page manuscrite 341.
de cette section ; elles se présentèrent en grand nombre et devinrent l'objet principal de leur sollicitude.

Pendant que la Congrégation des hommes achevait de se constituer, qu'elle s'accroissait à vue d'œil et qu'elle prenait tous les jours plus de consistance, la Congrégation des jeunes gens ne s'endormait pas ; malgré la précaution qu'on avait prise de ne lui enlever ses membres que successivement pour ne pas trop l'énerver, elle avait été un peu affoiblie par suite de la division. Tous les jours elle appelait bien quelques nouveaux frères à entrer dans la famille de Marie, mais malgré l'espèce d'urgence où elle se trouvait de faire des réceptions, elle se serait bien gardée de s'écarter le moins du monde des délais que la prudence exigeait. Elle n'oubliait pas qu'il fallait avant tout rester fidèle à la vieille maxime inscrite tout au long dans le règlement depuis l'origine de la Congrégation, qu'il fallait songer beaucoup moins à s'augmenter rapidement qu'à se composer solidement. Aussi les réceptions n'avaient encore pu suffire qu'à réparer une partie de ses pertes ; mais cet état de choses n'avait pas refroidi le zèle, et si le nombre des membres avait un peu diminué, l'ardeur qui les animait n'avait fait que s'accroître et apportait ainsi sous ce rapport une honorable compensation.

Plus que jamais, les sections s'efforçaient de diriger leurs travaux de manière à procurer la gloire de Dieu et le salut du prochain.

\section*{Sections}
% Mention marginale de la page manuscrite 342 ; déplacée avant la suite du même paragraphe dans l'édition.
Ce fut vers ce temps-là que la section de l'instruction apporta un notable changement à une de ses œuvres. Nous l'avons déjà vue, presque dès l'origine de la Congrégation, entreprendre
% Source : PDF original, vue 171, côté gauche, page manuscrite 342.
de former dans plusieurs paroisses de la ville des catéchismes de persévérance ; nous l'avons encore vue aussitôt après le retour de la paix reprendre ces catéchismes et leur donner plus d'extension ; aujourd'hui elle va faire plus encore. Sous le nom de société de St Louis de Gonzague elle les transforme en autant de petites Congrégations qui dirigées par les membres de la section reçoivent par cette impulsion comme une émanation de l'esprit qui anime la nôtre de laquelle elles demeurent une dépendance. Vous connaissez, Messieurs, parce que vous avez été les témoins du bien que ces sociétés ont opéré, et celui qu'elles opèrent encore tous les jours dans plusieurs paroisses de Lyon, notamment à St Bonaventure, à St Paul, à St Louis, mais par-dessus tout à St Polycarpe.

\section*{Section des prisons}
% Mention marginale de la page manuscrite 344, placée au début du développement p. 342.
Pendant que les travaux de la section de l'instruction donnaient des résultats aussi avantageux, la section des prisons et hôpitaux faisait aussi de son côté valoir le talent qui lui était confié. Elle était enfin parvenue grâce à l'intervention d'un membre de la Congrégation des hommes administrateur des prisons, à faire sa rentrée dans les prisons civiles, et elle y était protégée ; elle en profita pour y déployer plus de zèle que jamais. Elle se mit de nouveau à la disposition de MM. les aumôniers pour seconder de tout son pouvoir leurs travaux apostoliques.

Vers ce temps-là on donna une mission aux prisonniers de St Joseph ; la section fit tout ce qui dépendait d'elle pour en favoriser le succès. Ses membres se relayaient pour assister aux exercices ; ils parcouraient ensuite les chambres des prisonniers. Là ils employaient toute l'éloquence qu'une charité
% Source : PDF original, vue 171, côté droit, page manuscrite 343.
compatissante avait placée dans leurs cœurs, pour persuader à ces malheureux de ne pas laisser échapper une occasion aussi favorable de revenir enfin à ce Dieu de bonté qui les pressait de se convertir par des paroles pleines d'une tendresse toute paternelle que les Missionnaires leur adressaient chaque jour en son nom.

Ces touchantes exhortations ne demeurèrent pas sans fruits ; de leur côté, les Missionnaires s'étaient acquittés avec un zèle plein d'ardeur, de ces fonctions qu'ils remplissaient à la suite de ce bon Pasteur qui charge sur ses épaules la brebis égarée pour la ramener au bercail. Le souverain maître, père des Missionnaires, avait donné sa bénédiction à cette ligue toute sainte dont le but était de transformer des malheureux voués au mépris de la société qui les repousse, en des pécheurs pénitents devenus aux yeux de celui dont les pensées ne ressemblent pas aux pensées des hommes, ni ses jugements à leurs jugements, de véritables adorateurs en esprit et en vérité.

La Mission eut un succès complet ; mais le zèle de la section ne se contenta pas d'y avoir contribué ; elle voulut encore travailler à assurer le fruit de cette victoire remportée sur l'Enfer. Pour cela il fallait trouver un moyen de fixer les nouveaux convertis dans la pratique du bien, et les prémunir contre les dangers des mauvaises compagnies, des mauvais conseils et surtout des vieux souvenirs qui, lorsque ces hommes qui ont vieilli dans la carrière du crime se proposent d'abandonner la triste voie qu'ils suivent depuis si longtemps pour entrer enfin dans le chemin de la vertu, semblent se dresser devant eux comme des spectres pour leur
% Source : PDF original, vue 172, côté gauche, page manuscrite 344.
barrer le passage. Ils imaginèrent pour parvenir à ce résultat si désiré un projet dont l'exécution pouvait bien paraître chimérique, et cependant dont l'heureux succès est venu confirmer toutes leurs espérances ; ils organisèrent au milieu des prisonniers une Congrégation avec ses prières communes, ses fêtes, ses assemblées et ses divers exercices. Là les convertis continuèrent à s'édifier les uns les autres. Les bons conseils et les bons exemples qu'ils se donnaient réciproquement, ceux qu'ils recevaient de ceux de nos Congréganistes qui étaient chargés de les diriger, contribuèrent à les maintenir dans leurs bonnes dispositions. Cette pieuse institution fut dans le principe l'objet des railleries de quelques prisonniers. Elle a triomphé de cette petite persécution par la patience, et fini par conquérir l'estime générale et par être respectée même par les hommes les plus dépravés que renfermait la maison.

La seule chose qu'il n'était pas au pouvoir de la section de donner à cette Congrégation, c'était la fixité. On conçoit que dans une maison où ceux qui l'habitent sont sans cesse renouvelés, cette société n'a pu subsister que jusqu'au moment où ses principaux membres ont fait place à des nouveaux venus qui n'avaient point assisté à la mission, et qui par conséquent venaient tristement rendre à la prison dans laquelle ils étaient renfermés, l'aspect affligeant qu'on y remarque ordinairement.

La section des prisons profita de la liberté qui lui était rendue pour joindre une nouvelle œuvre à toutes celles qu'elle avait exercées précédemment. Celle-là paraît tellement pénible à la foiblesse de notre nature, qu'on ne peut pas douter qu'il n'y a que l'esprit de foi qui puisse donner la force de
% Source : PDF original, vue 172, côté droit, page manuscrite 345.
se dévouer à une œuvre pareille. Vous avez déjà pressenti que je veux parler ici du service immense de charité que les membres de la section des prisons et hôpitaux rendent aux condamnés à mort.

C'est le greffier de la cour d'assise qui d'après la loi devait annoncer aux malheureux condamnés que la cour de cassation a rejeté son pourvoi. Mais pour se débarrasser d'une corvée pénible, il se borne ordinairement à en prévenir le concierge de la prison. Avant que la section se fût chargée de cette œuvre le concierge désignait un des prisonniers pour aller porter au pauvre patient cette triste nouvelle. C'était assez ordinairement la fonction de celui qui était chargé de balayer la maison. Il avait pour mission d'étourdir le condamné en l'engageant à bien boire et à bien manger ; ce prisonnier ne demandait pas mieux car il en tirait sa part. On ne refusait rien au malheureux prêt à monter sur l'échafaud ni à son commensal. La dernière nuit se passait donc à se gorger de viandes et de vin, et lorsqu'on venait chercher la malheureuse victime on la trouvait ordinairement dans un état complet d'ivresse. C'est ainsi qu'elle était conduite au supplice, accompagnée par un vénérable prêtre dont les exhortations étaient rendues entièrement inutiles, par les dispositions et l'état d'anéantissement où se trouvait ce malheureux si tristement préparé pour paraître devant le souverain juge.

\section*{Section des prisons. Œuvre des condamnés à mort}
% Mention marginale de la page manuscrite 346, placée avant le début du récit p. 345.
Voici maintenant comment les choses se passent depuis que la section a apporté là sa généreuse assistance. Aussitôt que la section apprend qu'une condamnation à mort est prononcée, elle délègue deux de ses membres pour aller
% Source : PDF original, vue 173, côté gauche, page manuscrite 346.
visiter le condamné dans son cachot. Ceux-ci répètent fréquemment leurs visites pendant tout le délai que donne à l'exécution le pourvoi en cassation. Ce laps de temps est employé à essayer d'obtenir la confiance du malheureux qui doit mourir. Lorsque le rejet du pourvoi est arrivé, le concierge de la prison trop heureux d'avoir à sa disposition ce moyen de préparer la victime, se hâte de faire avertir nos Messieurs. Ceux-ci font à l'instant recommander le patient aux prières de tous les membres de la section, et courent s'enfermer dans son cachot pour passer avec lui cette dernière nuit. Leur présence à une heure indue, suffit ordinairement pour faire deviner à l'infortuné la fatale nouvelle. Dès ce moment tous leurs efforts sont employés à préparer cette pauvre âme à paraître devant Dieu ; ils implorent pour elle la clémence divine, par l'intercession de tous les saints, mais surtout par l'intercession de Marie le refuge des pécheurs ; ils offrent tour à tour à ses méditations, la justice de Dieu et la grandeur de ses miséricordes ; ils s'efforcent de lui inspirer tout à la fois un repentir sincère de ses péchés et une grande confiance dans celui qui a bien voulu mourir pour les expier ; ils prient avec le condamné et l'aident à préparer une confession bien sincère. Ce zèle si touchant est presque toujours couronné d'un succès complet. Le vénérable aumônier vient étendre sur ce pauvre pénitent cette main paternelle qui a reçu le pouvoir de lui fermer l'enfer et de lui ouvrir le ciel, et met ainsi le sceau à l'exercice de charité dont il est l'objet ; nos amis y persévèrent avec courage jusqu'au moment affreux où le bourreau se présente pour réclamer
% Source : PDF original, vue 173, côté droit, page manuscrite 347.
au nom de la justice humaine, la victime qu'il doit immoler sur l'échafaud ; alors nos Congréganistes avant de quitter pour la dernière fois celui pour lequel ils ont pris tant de soin et de peine, lui donnent en pleurant le baiser de paix, et pendant qu'accompagné par M. l'aumônier, il part pour aller consommer son sacrifice, eux vont se réfugier dans une église pour l'aider encore au moment suprême du secours de leurs prières.

Ce tableau sublime qu'offre ici une charité presque héroïque impose tellement à ceux qui sont le moins susceptibles de se laisser attendrir par des émotions de ce genre, qu'on a vu plusieurs fois les bourreaux et les gendarmes pénétrés d'admiration à la vue de ce ministère angélique exercé par des jeunes gens du monde au nom de la religion, demeurer immobiles à la porte du cachot et ne se décider à passer outre et à procéder aux tristes préparatifs de la réparation exigée par la justice humaine, que lorsque nos amis eux-mêmes, s'étant aperçus de cette cruelle arrivée, ont compris que le moment de se retirer était arrivé.

\section*{Marche de la Congrégation}
% Mention marginale de la page manuscrite 348, placée avant le bilan commencé p. 347.
Rien ne contribuait davantage à réchauffer le zèle des Congréganistes que des œuvres de ce genre. Au surplus ils se montraient généralement disposés à seconder la Congrégation de tout leur pouvoir ; elle commençait alors, appuyée sur les six sections, à marcher avec courage à la conquête des âmes ; non seulement elle se fortifiait en donnant comme un aliment à la bonne volonté de ses membres des œuvres nouvelles, mais l'esprit d'obéissance et de soumission qui animait tous les membres de l'une et l'autre Congrégation
% Source : PDF original, vue 174, côté gauche, page manuscrite 348.
était édifiant au-delà de toute expression. Le nombre des Congréganistes s'accroissait de jour en jour ; il ne tarda pas à s'élever de cent à cent vingt, et peu d'années après on en comptait cent soixante à cent soixante-dix. Les préfets, parfaitement unis, placés comme des frères dans la main du père de la Congrégation, leur centre commun, n'avaient à eux deux qu'une seule volonté ; cette douce harmonie favorisait merveilleusement la marche ascendante de la Congrégation ; plus nous avancions, plus nous avions à bénir la Providence de nous avoir inspiré l'heureuse pensée de cette division en deux Congrégations, qui avait déjà assuré à notre société des avantages beaucoup plus grands que ceux dont elle avait joui jusqu'alors et bien au-delà de tout ce qu'elle avait prévu. Mais notre bon Maître nous conduisait pas à pas ; nous verrons bientôt comment il voulait encore augmenter l'influence que nous pouvions exercer, afin qu'il nous fût possible d'accomplir les desseins qu'il avait sur nous. Déjà par suite de notre accroissement et de l'entrée d'un grand nombre de chrétiens notables dans la Congrégation des hommes, ce n'était plus seulement, comme dans les premières années, dans les confréries du St Sacrement que la Congrégation faisait sentir sa présence. Elle avait pénétré partout ; on la rencontrait dans les conseils de fabrique, dans les bureaux de bienfaisance, parmi les administrateurs des différents hospices, des prisons, et jusque dans le conseil municipal et dans les adjoints de la mairie. Si ses membres n'étaient pas assez nombreux dans ces différents corps pour y avoir la majorité, du moins ils pouvaient
% Source : PDF original, vue 174, côté droit, page manuscrite 349 (début).
y influer assez sur l'opinion de leurs collègues pour rendre ces différentes administrations favorables aux progrès de la religion dans notre ville. Il est facile d'apprécier ce qu'y pouvait faire dans l'intérêt du bien l'action d'hommes unis par des liens intimes, et plus encore par l'esprit de foi qui dirigeait tous leurs cœurs dans une pensée commune. C'était déjà un bien grand succès que la Congrégation avait dès lors obtenu ; mais ne nous lassons pas de le redire, cette union toute spirituelle n'a jamais été employée pour des intérêts temporels ; notre royaume n'est pas de ce monde.

Aussitôt après la division, la Congrégation sollicita et obtint du souverain Pontife un rescrit qui étendait la grâce des indulgences aux membres de la Congrégation des hommes. Ainsi des indulgences communes devinrent encore un nouveau lien entre les deux Congrégations.
